\section{Comparison of methods}
\label{sec:comparison}

Figure~\ref{fig:taxonomy} shows how the methods are grouped, and Table~\ref{tab:comparison} lists
them side by side. The question the table answers: for a given set of available measurements,
which methods exist, which model do they assume, and where is each described?

\begin{figure}[t]
  \centering
  % Method families and what each needs as input. Included by ../sections/comparison.tex
% with % Method families and what each needs as input. Included by ../sections/comparison.tex
% with % Method families and what each needs as input. Included by ../sections/comparison.tex
% with \input{figs/taxonomy}. Keep labels to a few words; explain in the surrounding prose.
\begin{tikzpicture}[
    fam/.style  = {draw, rounded corners=2pt, align=center, text width=24mm,
                   minimum height=9mm, font=\small, fill=blue!8},
    need/.style = {font=\footnotesize, align=center, text=black!70},
]
  \node[fam]                  (a) {Family A};
  \node[fam, right=6mm of a]  (b) {Family B};
  \node[need, below=3mm of a] {needs: \todo{input}};
  \node[need, below=3mm of b] {needs: \todo{input}};
\end{tikzpicture}
. Keep labels to a few words; explain in the surrounding prose.
\begin{tikzpicture}[
    fam/.style  = {draw, rounded corners=2pt, align=center, text width=24mm,
                   minimum height=9mm, font=\small, fill=blue!8},
    need/.style = {font=\footnotesize, align=center, text=black!70},
]
  \node[fam]                  (a) {Family A};
  \node[fam, right=6mm of a]  (b) {Family B};
  \node[need, below=3mm of a] {needs: \todo{input}};
  \node[need, below=3mm of b] {needs: \todo{input}};
\end{tikzpicture}
. Keep labels to a few words; explain in the surrounding prose.
\begin{tikzpicture}[
    fam/.style  = {draw, rounded corners=2pt, align=center, text width=24mm,
                   minimum height=9mm, font=\small, fill=blue!8},
    need/.style = {font=\footnotesize, align=center, text=black!70},
]
  \node[fam]                  (a) {Family A};
  \node[fam, right=6mm of a]  (b) {Family B};
  \node[need, below=3mm of a] {needs: \todo{input}};
  \node[need, below=3mm of b] {needs: \todo{input}};
\end{tikzpicture}

  \caption{Method families and what each needs as input. \todo{Fill in.}}
  \label{fig:taxonomy}
\end{figure}

% Wide table across both columns. Add a row per method (or per group of methods that share
% all four descriptive columns). Never invent an entry: every row needs a reference.
\begin{table*}[t]
\caption{State-of-the-art summary. \todo{Fill in from the literature notes.}}
\label{tab:comparison}
\small
\begin{tabularx}{\textwidth}{@{}L L L L l@{}}
\toprule
\textbf{Measurements used} & \textbf{Estimated quantity} & \textbf{Model used} &
\textbf{Methodology} & \textbf{Ref.}\\
\midrule
\todo{signals} & \todo{quantity} & \todo{model} & \todo{technique} & \cite{singh2019}\\
\addlinespace
\todo{signals} & \todo{quantity} & \todo{model} & \todo{technique} & \todo{ref}\\
\bottomrule
\end{tabularx}
\end{table*}

\subsection{Family A}
\todo{One subsection per family in the taxonomy: the idea in plain words, what it needs, where it
works and where it fails, and which table rows belong to it.}

\subsection{Family B}
\todo{As above.}
